\documentclass[10pt,a4paper]{amsart}
\usepackage[margin=19mm]{geometry}
\usepackage{amsmath,amssymb,mathtools}
\usepackage[scr=rsfs,cal=euler]{mathalfa}
\usepackage{xfrac}
\usepackage[unicode,hidelinks,bookmarks=false]{hyperref}
\newcommand{\ie}{i.e.\ }
\newcommand{\cf}{cf.\ }
\newcommand{\resp}{respectively\ }
\newcommand{\Subsection}{\subsection}
\providecommand{\nobreakdash}{\nobreak\hbox{-}}
\setlength{\emergencystretch}{2em}
\multlinegap0pt
\newcommand{\cld}{\text{Ch}}
\newcommand{\ind}{\textit{ind}}
\usepackage{amssymb}
\newtheorem{remark}{Remark}
\newtheorem{prop}{Proposition}
\newtheorem{lemma}{Lemma}
\newcommand{\marked}{\mu}
\newcommand{\pair}[1]{\langle #1\rangle}
\title{Matching Rules for Substitution and Hierarchical Tilings for any Substitution with Finite Local Complexity}
\author{Nikolay Vereshchagin}
\date{Source: arXiv:2606.25005v2 (CC BY 4.0)}
\begin{document}
\maketitle

\noindent\emph{Source and licence.} Matching Rules for Substitution and Hierarchical Tilings for any Substitution with Finite Local Complexity. Version arXiv:2606.25005v2 (CC BY 4.0), distributed by arXiv under the Creative Commons Attribution 4.0 International licence, http://creativecommons.org/licenses/by/4.0/, as recorded in the arXiv licence metadata for this exact version and shown on the abstract page https://arxiv.org/abs/2606.25005v2. Full licence notice: \texttt{licenses/arxiv-2606.25005-CC-BY-4.0.txt}.\par\smallskip

\noindent\textit{Editorial scope.} Counted unit: the article's prop th2 (source0.tex lines 1087--1090). The article's complete original proof of it is delivered with it (source0.tex lines 1091--1537, one \texttt{proof} environment of 447 lines). No mathematics is written, repaired or re-derived: the statement and the proof are the article's own lines, in the article's own notation, and no retained line is substituted or re-flowed.  Retained uncounted as the source-local context the proof needs: lines 560--560; lines 857--861; lines 950--967.  One declared adaptation: exactly 3 ancillary strings inside the retained spans are omitted (image-inclusion commands and, where the source repeats a label, the repeated label definitions) and no other line differs from the source; the figure environments, with their placement, captions and labels, are kept, and the package records the omitted strings in fidelity receipts bound to the affected bodies.  No retained line contains a citation, so the wrapper carries no bibliography and no bibliography entry.  The retained lines contain the article's own diagram code, which the wrapper typesets with the standard tikz packages; no image file is delivered and the package declares no asset dependency.  Editorial formatting only, in addition: an AMS \texttt{amsart} preamble that restates the article's own declarations and macros used by the retained lines, namely the environments \texttt{remark}, \texttt{prop}, \texttt{lemma} on separate counters, together with the standard packages those lines need.  The retained environments print in this document as Remark 1, Proposition 1, Lemma 1 in order of appearance, whereas the article prints the same items with the numbering of its own full document.  The complete native source of the article as distributed in the arXiv e-print for this exact version is bundled unchanged as \texttt{sources/203574/source0.tex}.
\subsubsection{Marked sides of prototiles from $P'$}\label{smarked}

\begin{remark}\label{rem-parent}
A tile may have several parents or no parents at all. In the former case, all its parents
are of the same form. If $B$ is a boundary tile, then, moreover, all its parents
are of the same type. All tiles that have a parent are normal.
\end{remark}

\renewcommand{\labelenumi}{S\theenumi:}
\begin{enumerate}
\item Every legal tile has a parent and hence is  normal --- its parent is 
the second term $A_1$ of the
ancestor sequence that confirms legality.
\label{s1}
\item Every child of every legal tile is also legal.\label{s3}
\item The parent indices on different sides of a legal tile
match. \label{s2}
\item
If a legal tile has two sides on the path $P_i(a)$
then the copies of  $\ind_A(i,a)$
on those sides 
match.\label{s6}
\item
If two sides $a,b$ of the same tile share a common vertex $V$, then the crowns for $V$
in the crown indices of these sides match.\label{s7}
\end{enumerate}

\begin{prop} \label{th2}
  Every proper tiling $\mathcal T$ with legal tiles has a 
proper $\sigma''$-composition $\mathcal T'$ consisting of legal tiles.
 \end{prop}

\begin{proof}
Recall that a central type $t$ is called \emph{cyclic} if $t$ is the central type in the macrotile $\sigma' t$, see Fig.~\ref{f34}.
\begin{figure}
\begin{center}
\hskip 2.5cm 
\end{center}
\caption{The central type 1 is acyclic, and the central type 2 is cyclic.}\label{f34}
\end{figure}
The key lemma in the proof is the following 
\begin{lemma}[Composition Lemma]\label{l9}
Let $\mathcal M$ be a finite proper tiling 
that is equal to a $\sigma'$-macrotile 
$\sigma' \alpha$  provided we ignore all indices except the sixth one (the one written in the
middle of a tile, which specifies its type).
Assume that all tiles in $\mathcal M$ are legal, except possibly for the  central tile $A$. 
    Then the following hold: 
    \begin{enumerate}
    \item[(a)] All boundary sides of  $\mathcal M$ have the same parent index $t$, which is of the form $\alpha$.
  \item[(b)] There is a normal tile $D$ of type $t$ with $\sigma''D\ni\mathcal M$.
      
    \item[(c)] If the central tile $A$ is legal and  $t$ is a non-central or cyclic type, then $D$ is legal as well.
    \end{enumerate}
  \end{lemma}
  \begin{proof}
    (a) Due to property~S\ref{s2} and since $\mathcal M$ is proper, all the boundary sides of  $\mathcal M$ have the same parent index $t$. To prove that $t$ is of the form $\alpha$, 
it suffices to find a legal type-$t$ tile of the form $\alpha$.
   Such a tile is any legal parent $C$ of any border tile $B$ in $\mathcal M$,
   see Fig.~\ref{f8}.
    \begin{figure}
      \begin{center}
        
      \end{center}
      \caption{On the left is a proper tiling $\mathcal M\in \sigma'' D$ with $\pi(\mathcal M)=\sigma'\alpha$ 
and with the  central tile $A$. 
        On the right is the legal parent $C$ of $B$ and a decomposition of $C$. 
        The type of $C$
        and the parent indices in its decomposition and in the tiling $\mathcal M$ are  denoted
        by $t$.
}\label{f8}
    \end{figure}
   The form of $C$ is $\alpha$, since $C$ has a child in $\mathcal M$ and $\pi(\mathcal M)=\sigma'\alpha$. And the type of $C$ is $t$, since so is
 the parent index of $B$.
    

(b) Let $D$ be the following tile, whose indices are determined
by $t$ and third  indices on ports of  $\mathcal M$:
\begin{itemize}
\item the sixth index of $D$ is $t$,
\item the $i$th index of the side $d$ of $D$ is equal to the third index on the $i$th 
port on the macroside $\sigma'd$ of $\mathcal M$.
\end{itemize}


Let us show that this tile $D$ is normal. Let $\ind_D(i,d)$ be a local
index of $D$. 
By construction, 
$$
\ind_D(i,d)=\ind_B(3,b)
$$
where $b$ is the $i$th port
on the macroside $\sigma'd$ of the macrotile $\mathcal M$ and 
$B$ the tile from $\mathcal M$ containing this port (see Fig.~\ref{f8}). 
By assumption,
tile $B$ is legal, and therefore has some legal parent $C$.
The type of $C$ is $t$, since that is the parent index of $B$. Indeed, the tile $B$ contains a
port and hence has an outer side; not all sides of $B$ are outer, so $B$ has a boundary side as
well. By item (a) the parent index on that side equals $t$, and by the requirement on the fourth
index of boundary sides it equals the sixth index of $C$, that is, the type of $C$.
Furthermore,
$$
\ind_B(3,b)=\ind_C(i,d)
$$
by the definition of the substitution $\sigma''$.
Therefore, $\ind_D(i,d)=\ind_C(i,d)$,
and therefore $\ind_D(i,d)$ takes the value allowed for normal tiles of type $t$.
Note that in this argument it matters only that the value of a local index for a normal 
tile is a function of $t$ and $d$.

Let us show
that $\mathcal M$ is a $\sigma''$-decomposition of $D$.
First, note that $\mathcal M$ satisfies the Crown Consistency Property,
since all the tiles that have vertices on the border of $\mathcal M$ are  legal and 
thus satisfy Property~S\ref{s7}. Therefore,
it suffices to prove that
for all types $s$ in $\sigma'\alpha$, the tile $\mathcal M_s$
has indices satisfying the requirements for indices of tiles
$\mathcal M_s$ for $\mathcal M\in\sigma''D$.

First, assume that tiles of type $s$ do not have borrowed indices; in particular, type $s$ is neither central nor boundary.
Then, each tile of the form $\alpha$
has only one child of type $s$, and for different tiles of the form $\alpha$,
these children coincide.
Since the tile $\mathcal M_s$ is legal, it has a parent and that parent must be of the form $\alpha$.
Therefore, $\mathcal M_s$ is a child of $D$.

Now suppose that type $s$ has borrowed indices but is not central. Let $\ind_{\mathcal M_s}(j,b)$ be any index of the tile $\mathcal M_s$.
We need to prove that this index satisfies the requirements for indices
of tiles from macrotiles in $\sigma'' D$.
By assumption, the tile $\mathcal M_s$ is legal, and therefore
has a legal parent $C$. Therefore, $\ind_{\mathcal M_s}(j,b)$
satisfies the requirements for macrotiles in $\sigma'' C$. If
$\ind_{\mathcal M_s}(j,b)$ is not borrowed, then these requirements for
$\sigma'' C$ and $\sigma'' D$ are the same, so $\ind_{\mathcal M_s}(j,b)$
also satisfies the requirements for macrotiles from $\sigma'' D$.

If
$\ind_{\mathcal M_s}(j,b)$ is borrowed, then the argument is a little more complicated.
Recall that there are three sorts of borrowed indices: the parent index on a boundary side, the
third index on a non-special side belonging to a path, and the first four indices on a special
side that is a print. We consider these three cases in turn.

\emph{The parent index} ($j=4$ and $b$ is a boundary side). Then $\ind_{\mathcal M_s}(j,b)$
must be equal to the type of tile $D$. This is indeed the case
by the construction of tile $D$.

\emph{The third index on a path.} Assume that $b$ is not special and $j=3$, so that
$b\in P_i(d)$ for some $i\le5$ and some side $d$ of tile $D$.
In this case, we need to prove that $\ind_{\mathcal M_s}(3,b)=\ind_D(i,d)$.
Indeed, all sides of the path $P_i(d)$ in the macrotile $\mathcal M$, except the print at which
the path ends, have the same third index. For two sides of the path that belong to the same tile
this follows from Property~S\ref{s6}: both their third indices are copies of one and the same
index of a legal parent of that tile. For two sides of the path that are shared by neighbouring
tiles this follows from the properness of $\mathcal M$.
In particular, $\ind_{\mathcal M_s}(3,b)$ is  equal to the third index of the $i$th port of the macrotile $\mathcal M$
on the macroside
$\sigma' d$. By the construction of the tile $D$, the latter is equal to $\ind_D(i,d)$.

\emph{The borrowed indices on a print.} Assume finally that $b$ is special and is the print of a
marked side $d$ of the form $\alpha$; then $j\le 4$, and what we have to prove is that the color
of $b$ in $\mathcal M_s$ equals
$\pair{\ind_D(1,d),\ind_D(2,d),\ind_D(3,d),\ind_D(4,d),I_s(b)}$, as $\sigma''D$ requires. Since
the type $s$ is not central, $s$ is the small quasi-square that shares the print $b$ with the
central type. The tile $\mathcal M_s$ is legal, hence it has a legal parent $C$, and the form of
$C$ is $\alpha$, because $C$ has a child of type $s$. Consequently $C$ and $D$ have the same
prints and the same paths, and the requirements of $\sigma''C$ tell us that the color of $b$ in
$\mathcal M_s$ is $\pair{\ind_C(1,d),\dots,\ind_C(4,d),I_s(b)}$. So it suffices to show that
$\ind_C(j,d)=\ind_D(j,d)$ for every $j\le4$.

Consider the side $c$ of the path $P_j(d)$ that precedes the print $b$. The path has at least
three sides, since no tile of the macrotile has both an outer side and a side of the central tile.
By the choice of the paths, $c$ is a non-special side of the small quasi-square $s$ itself, so the
requirements of $\sigma''C$ give $\ind_{\mathcal M_s}(3,c)=\ind_C(j,d)$. On the other hand, by the
previous case the third index is the same on all sides of $P_j(d)$ except the print, so
$\ind_{\mathcal M_s}(3,c)$ equals the third index of the $j$th port of $\mathcal M$ on the
macroside $\sigma'd$, which is $\ind_D(j,d)$ by the construction of $D$. Hence
$\ind_C(j,d)=\ind_D(j,d)$, which is what we needed.




   It remains to show that $A$, the central tile in $\mathcal M$, is in  $\cld_s(D)$, where
   $s$ is the central type in $\mathcal M$. 
     We have to prove that on each side $a$ of $A$ all
    indices are as required by the definition of $\cld_s(D)$. Choose any such side $a$.
    On that side, $A$ has the same indices as its neighbour $B$ on side $a$,
    since the tiling $\mathcal{M}$ is proper.
    And $B$  has the indices required for the child of $D$ of its type.
    By definition of $\sigma''$,  the indices of adjacent children on shared interior sides coincide.
    Therefore, the tile $A$ has the required indices on the side $a$. Note that here we use the
    symmetry of the requirement on prints: if $a$ is a print, then $\sigma''D$ prescribes on $a$
    the same borrowed tuple to the central type and to the small quasi-square sharing $a$ with it.

    (c) 
First, consider the case where $t$ is a cyclic type. We
claim that in this case $D$ and $A$ coincide, hence $D$ is legal.
Indeed, by construction, $D$ is of type $t$, thus the central child of $D$,
that is, the tile $A$, is also of type $t$. Both tiles $D,A$ are normal, therefore 
they have the same local indices. Finally, since $A$ is 
the central child of $D$,
its nonlocal indices copy those of $D$.

We now turn to the case where $t$ is a non-central type.

Consider any chain of tiles
$$
A_0=A,A_1,A_2,\dots\in P'',
$$
that proves the legality of $A$.
We claim that at least one tile in this sequence has a non-central type.
Indeed, otherwise, the first index on all marked sides of $A$ would be zero.
But we know this is not so, since $A$ inherits its first index from $D$,
which is not zero on at least one side; recall that $D$ is a normal non-central tile
and hence has an inner marked non-special side (prints are defined for marked sides only, and both
kinds of non-central types --- original tiles and small quasi-squares --- do have such a side).

Now consider the shortest chain
$$
A_0=A,A_1,A_2,\dots,A_n\in P'',
$$
in which each tile is a central child of the next one, and the last tile $A_n$
is of a non-central type and is legal. 
The central child of a central tile has the same form as the tile itself, so their central children
have the same type. Therefore, the types of all the tiles $A_0,\dots,A_{n-2}$ coincide; denote this common type by $s_0$.
Since $A_0$ is the central child of $A_1$ and has type $s_0$ too, the central child of a type-$s_0$ tile again has type $s_0$,
that is, $s_0$ is cyclic. By the cyclic case treated above, a normal tile of a cyclic type coincides with its central child;
hence the tiles $A_0,\dots,A_{n-2}$ all coincide.
Hence we can remove from the chain all the tiles $A_1,\dots,A_{n-2}$, which implies that
$n\in\{1,2\}$.

Let us show that in fact $n=1$. Assume that $n=2$. The tiles $A_1$ and $D$ are both parents
of $A$, hence by Remark~\ref{rem-parent} they have the same form, namely the form $\alpha$ of the
type $t$. On the other hand, if $n=2$, then $A_1$ is the central child of $A_2$, and the central
tile of a macrotile is always a large quasi-square, so the form of $A_1$ is a large quasi-square.
But the form of the non-central type $t$ is a small quasi-square or an original prototile, and
never a large quasi-square --- a contradiction. Thus $n=1$; let $s$ denote the type of the legal
tile $A_1$.

We claim that  $s=t$. To prove the claim, 
denote by $\Phi(A)$ the multiset consisting of all types $r$ such that 
there is a type $u$ such that  $\pair{r,u}$ or $\pair{u,r}$ is a first index of a \emph{marked} side of $A$.
%of components of nonzero first indices on the sides of tile $A$.
The cardinality $|\Phi(A)|$ is 2 times the number $\marked(A)$ of marked inner sides of $A$ provided
$A$ has no special sides. Indeed, the first index on every non-special marked inner side is the identity
index, which contributes 2 to $|\Phi(A)|$. If $A$ has a special side (which can happen only if $A$ is a small quasi-square),
then $|\Phi(A)|$ may be $2\marked(A)-2$, since the first index on the special side of $A$ can be arbitrary.

We have  $\Phi(A_1)=\Phi(D)$; in fact the first four indices of $A_1$ and of $D$ coincide side by
side. Indeed, $A_1$ and $D$ are parents of one and the same tile $A$, so by
Remark~\ref{rem-parent} they are of the same form $\alpha$, and therefore one and the same print
map $a\mapsto b(a)$ is used for both of them (prints are chosen for a prototile, so they depend
only on the form). All marked sides of the central child $A$ are special, and $b(a)$ is a print, so
the definition of $\sigma''$ gives, for every marked side $a$ of $\alpha$ and every $j\le4$,
$$
\ind_{A_1}(j,a)=\ind_A(j,b(a))=\ind_D(j,a).
$$
In particular, the multisets $\Phi(A_1)$ and $\Phi(D)$ coincide.

Let us stress that the last argument compares $A_1$ and $D$, and \emph{not} $A$ and $D$:
the equality $\Phi(A)=\Phi(D)$ is in general false. The point is that all $4N$ marked sides of the
central child $A$ are special, while the number of prints equals the number of marked sides of
$\alpha$, which is smaller than $4N$ when $\alpha$ is not a large quasi-square; on the special
sides of $A$ that are not prints the first index is the identity index, and it contributes to
$\Phi(A)$ pairs that have nothing to do with $D$.

Thus it suffices to prove the following

\begin{lemma}\label{l8}
Assume that $A,B$ are normal non-central tiles of different types $s,t$. Then $\Phi(A)\ne \Phi(B)$. 
%Then $A$ and $B$ are of the same type.
\end{lemma}
\begin{proof}
By way of contradiction assume that $\Phi(A)= \Phi(B)$. 

Without loss of generality we may assume that the number of marked sides of each small quasi-square 
is much larger than the number of sides of any of the original tiles.
Indeed,  let $\Gamma$
denote the largest side length of the prototiles from $P$.
Then the number of marked sides of any small quasi-square is at least roughly $3R/\Gamma$; recall that $R$ denotes
the side length of the grid square. % (disregarding boundary effects, which do not affect the estimate much). 
Here the factor 3 comes from the fact that 
the side of the small quasi-square that is adjacent to the large quasi-square has only a few marked sides.

Recall that we may choose $R$ arbitrarily large, since
the power $\tau^k$ of the original substitution can be chosen arbitrarily
large. Indeed, it is enough that a grid of size $(N+2)\times (N+2)$ and $4\cdot4\cdot N$ disjoint paths fit inside the
$\tau^k$-macrotiles, where $N$ depends only on the original set and on $R$. Thus
we can first choose $R$, then choose $N$, and only afterwards choose a
sufficiently large $k$.

Hence $\Phi(A)=\Phi(B)$ implies that both $A,B$ are either original tiles having the same number of inner sides,
or small quasi-squares, whose number of marked sides differs by at most 1.

Note also that in both cases the types $s$ and $t$ must lie in one and the same macrotile, since
distinct macrotiles are disjoint. Indeed, if $A$ is an original tile, then it has no special sides,
so every member of $\Phi(A)$ is a component of an identity index of a side of $s$ and hence is a
type of the macrotile containing $s$; the same holds for $B$ and $t$, and $\Phi(A)=\Phi(B)$ is
nonempty. If $A$ and $B$ are small quasi-squares, then each of them has exactly one special side,
so all members of $\Phi(A)$ except at most two are types of the macrotile of $s$, and likewise for
$B$; as the number of marked sides of a small quasi-square is large, this again forces the two
macrotiles to coincide. So in what follows $s$ and $t$ are types of one macrotile $\sigma'\alpha$.

Assume first that $A,B$ are original tiles with $l$ inner sides,
so that $|\Phi(A)|=|\Phi(B)|=2l$. By the definition of identity index,
 $\Phi(A)$ includes $s$ with multiplicity at least $l$, while $\Phi(B)$ includes $t$ with multiplicity at least $l$.
Thus $s\ne t$ implies that all identity indices on inner sides of $A,B$ are $\pair{s,t}$ or $\pair{t,s}$.
In other words, types $s,t$ share all inner sides, which is impossible provided $\tau^k$-macrotiles
are large enough. Indeed, in that case every side on the boundary of the region $s\cup t$ is an
outer side of the macrotile, so $s\cup t$ is the entire macrotile, that is, the macrotile consists
of two tiles only.

Now assume that $A,B$ are small quasi-squares. %, whose number of marked sides differ by at most 1.
WLOG assume that $A$ is to the north of the central tile. Then $A$ has at least roughly $R/\Gamma$
northern sides. All of them except a constant number are shared with original tiles $C$
such that $B$ does not share any side with $C$. Each such side contributes to $\Phi(A)$ the type of
the corresponding tile $C$, and that type can occur in $\Phi(B)$ only among the at most two members
of $\Phi(B)$ that come from the special side of $B$, whose first index may be arbitrary. Thus
$\Phi(A)$ has approximately $R/\Gamma$
members that are outside $\Phi(B)$. As we have seen, we can make $R$ arbitrarily large, which implies a contradiction.
\end{proof}


%First consider the case where the forms of both tiles $A_n$ and $D$ are original prototiles,
%that is, belong to $P$.  

%Then the first indices on all sides of these tiles are identity indices, and
%on every inner side of $A_n$ one of the components of the identity index equals $s$. Likewise  
 %for $D$ one of the components of the identity index on inner sides equals $t$. 
%Since macrotiles are large enough, the type $s$ has at least one inner side that is not a side of the type $t$, or vice versa. 
%Assume, say, the first case.
%Then $\Phi(A_n)$ contains a pair of the form $\pair{s, r}$ or $\pair{ r,s}$ with $r\ne t$,
%while $\Phi(D)$ contains no such pair. Contradiction.

%In the remaining case, one of the tiles $A_n$ and $D$ is a small quasi-square, while the other 
%either belongs to $P$, or is also a small quasi-square.
%In this case, the argument is slightly more involved.
%Suppose, say, $A_n$ is a small quasi-square. 
%%Then $A_n$ may have two sides belonging to a path $P_1$,
%%and the first index on them may be arbitrary. In particular, this index need not be of the form 
%%$\pair{s, r}$ or $\pair{ r,s}$,
%%and conversely it may have the form $\pair{t, r}$ or $\pair{ r,t}$.
%%Furthermore, if $s$ and $t$ are neighboring small quasi-squares, then one of the indices
%%in $\Phi(A_n)$ or $\Phi(D)$ equals $\pair{s,t}$ or $\pair{ t,s}$. However, w
%Without loss of generality we may assume
%that every small quasi-square is adjacent to at least $K_0$ different original tiles along non-special sides,
%where $K_0$ is the maximal number of sides of original tiles. 
%Moreover, for any two distinct small quasi-squares, none of these $K_0$ tiles is adjacent to any other small quasi-square. 

%Indeed, let us call two original  tiles $A,B$ \emph{neighbours} if they have
%sides $a,b$ bordering the same original tile $C$. The distance
%between neighbouring tiles does not exceed the maximal perimeter of the
%prototiles from $P$, denoted in what follows by $\Gamma$. Let $\gamma$
%denote the smallest side length of the prototiles from $P$. And let $S$
%be some small quasi-square lying, say, to the north of the large
%quasi-square. It follows from the above that at most $\Gamma/\gamma$
%original tiles $A\subset S$ can be neighbours of some tile $C$ from the
%small quasi-square situated to the right of $S$. The same holds for the
%small quasi-squares to the left. On the other hand, the total number of
%neighbours of $S$ along its north side is at least $R/\Gamma$, where $R$
%is the side length of the grid square (disregarding boundary effects, which
%do not affect the estimate much). Therefore, if $R/\Gamma\gg K_0$, the
%claim holds. Now recall that we may choose $R$ arbitrarily large, since
%the power $\tau^k$ of the original substitution can be chosen arbitrarily
%large. Indeed, it is enough that a grid of size $N\times N$ fit inside the
%macrotiles, where $N$ depends only on the original set and on $R$. Thus
%we can first choose $R$, then choose $N$, and only afterwards choose a
%sufficiently large $k$.


%In this case
%$\Phi(A_n)$ contains $K_0$ indices with one component equal to $s$, and at least one of them does not belong to $\Phi(D)$.
%Contradiction.





%So we know that  $s=t$. 
Let us continue the proof of item (c). 
We showed that  tile $A$ has a legal parent $A_1$ of type $t$
and a normal parent $D$ also of type $t$.
We claim that $A_1=D$ and hence $D$ is legal.
As we have shown above, the tiles $A_1$ and $D$
have identical first four indices on all marked sides.
Their fifth indices on marked sides coincide as well, since $A_1$ and $D$ are normal tiles of the same type.
On non-marked sides all their indices are local and hence coincide for the same reason.
Therefore $A_1=D$, and so $D$ is legal. Lemma~\ref{l9} (Composition lemma) is proved.
  \end{proof}

Let us continue the proof of Proposition~\ref{th2}.
Let $\mathcal T$ be a proper tiling with legal tiles. 
Identity indices guarantee that it can be partitioned into $\sigma'$-macrotiles. For each of these macrotiles $\mathcal{M}$ there is a prototile $\alpha$ with $\pi(\mathcal{M})=\sigma'\alpha$.
Replace each macrotile $\mathcal M$ of the original tiling by the tile $D$ existing by 
the Composition Lemma.
  We obtain a $\sigma''$-composition $\mathcal T'$ of $\mathcal T$.
 Alignment indices on outer sides ensure that the composed tiling is  side-to-side.
Since $\mathcal T$ is proper, all indices of its tiles on adjacent ports match.
And since all indices on sides of composed tiles are borrowed from ports of $\mathcal T$,
the composed tiling $\mathcal T'$ is proper. 

  It remains to prove that each tile $D\in\mathcal T'$  is legal.
  By 
the Composition Lemma, $D$ is legal unless $D$ is of central
  acyclic type.
  We have to prove that  $D$ is legal even in that case.
  
Since all composed tiles are normal, $\mathcal T'$ can be again split into macrotiles.
Consider a macrotile $\mathcal M'\subset\mathcal T'$  containing a tile $D$ of a central acyclic  type $t$.
  Denote by $r$ the parent indices of tiles in $\mathcal M'$. By the Composition Lemma, the tiling $\mathcal M'$ has a composition 
$C$, which is a normal tile of type $r$,
see Fig.~\ref{f40}.
  \begin{figure}
    \begin{center}
      
    \end{center}
    \caption{Tiles $C,D,A$ have  types $r,t,s$, respectively, where $r$ is a non-central type and $t$
is a central acyclic type.
      The macrotile $\mathcal M'$ is a decomposition of $C$, and the macrotile $\mathcal M$ is a decomposition of $D$.
      The figure illustrates the proof that  $D$ is legal.}\label{f40}
  \end{figure}
   
Since the central child of $r$ is the acyclic type $t$,
the type $r$ is not central. Indeed, suppose $r$ were central. The form of a central type is a large quasi-square,
which is a fixed point of passing to the central child; hence $r$ and its central child $t$ have the same form.
Since the central child of a type depends only on its form, $t$ and $r$ have the same central child, namely $t$.
That is, the central child of $t$ is $t$ itself, so $t$ is cyclic --- contradicting the assumption that $t$ is acyclic.

Since $A$ is legal and is a grandchild of the normal tile $C$ of a non-central type, there exists a sequence of legal tiles
$$
A_0=A,A_1,A_2,\dots,A_n,
$$
in which each tile is a central child of the next one, the last tile is of a non-central type, and $n\in\{1,2\}$.
This is proven in exactly the same way as in the proof of item (c) of the Composition Lemma. (Only
the bound $n\in\{1,2\}$ is borrowed from there; the argument that excluded $n=2$ in item~(c) used
the non-centrality of the type of $D$, which is not the case now.)


Note that $A_1$ and $D$
are of the same form, as they have a common
central child $A$. Therefore $A_1$ is a central tile, hence $n=2$.
So, we have a chain of legal tiles $A=A_0,A_1,A_2$, in which each tile is a central child of the next one,
and $A_2$ is not central.

We claim that the tiles $A_1$ and $D$ have the same type. Exactly as in the proof of item~(c) of
the Composition Lemma, the tiles $A_1$ and $D$ are parents of one and the same tile $A$, hence
they are of the same form, hence the same print map $a\mapsto b(a)$ serves both of them, and
therefore
$$
\ind_{A_1}(j,a)=\ind_A(j,b(a))=\ind_D(j,a)
$$
for every marked side $a$ of their common form and every $j\le4$; in particular
$\Phi(A_1)=\Phi(D)$. This time, however, the tiles $A_1$ and $D$ are central, so
Lemma~\ref{l8} is not applicable to them. Instead we compare the multiplicities
of the type $t$ in these two multisets.

Being central tiles, both $A_1$ and $D$ have exactly $4N$ marked sides, and all of them are
special. The prints of the sides of $C$ are at most $Q$ of the marked sides of $D$, because the
form of the non-central type $r$ has at most $Q$ marked sides. On each of the remaining at least
$4N-Q$ special sides of $D$ the first index is the identity index of the corresponding side of the
central type $t$, and one of its two components equals $t$. Hence $t$ occurs in $\Phi(D)$ at least
$4N-Q$ times.

Assume now that the type $t'$ of $A_1$ differs from $t$, and count the occurrences of $t$ in
$\Phi(A_1)$. On a special side of $A_1$ that is not a print, the first index is the identity index
of a side of the central type $t'$; its components are $t'$ and the type of the small quasi-square
adjacent to that side, so neither of them is $t$ --- indeed $t'\ne t$ by the assumption, and $t$ is
a central type and hence is not a small quasi-square. The remaining marked sides of $A_1$ are
prints, and there are at most $Q$ of them; each side contributes at most two members to
$\Phi(A_1)$. So $t$ occurs in $\Phi(A_1)$ at most $2Q$ times. As $4N>3Q$ by the choice of $N$ in
Section~\ref{smarked}, we get $4N-Q>2Q$, which contradicts $\Phi(A_1)=\Phi(D)$. Therefore $t'=t$,
that is, $A_1$ and $D$ are of the same type.

Since  nonlocal indices of $A_1$ and $D$ coincide (they are inherited by their common child  $A$),
and their local indices coincide as well, being those of normal tiles of one and the same type,
the tiles $A_1$ and $D$ coincide, hence $D$ is legal.

Proposition~\ref{th2} is proved.
\end{proof}

\end{document}
